\subsection{群胚的轨道}

设 G 为群胚。G 的底箭图的连通分支称为 G 的轨道。G 的轨道所成的集记作 \(\mathrm{Orb}(\mathrm{G})\) ；若 \(\varphi: G \to G'\) 是群胚态射，则映射 \(\pi_{0}(\varphi)\) （由 \(\varphi\) 过渡到伴随图的连通分支而得的）通常记作 \(\mathrm{Orb}(\varphi)\) ，称为由 \(\varphi\) 过渡到轨道而得的映射。

恰有一个轨道的群胚称为传递的。若 G 是传递群胚，则诸群 \(G_{a}\) （对 \(a \in \text{Som}(G)\) ）彼此同构。若群胚 G 是传递的，且诸迷向群 \(G_{a}\) 都缩为各自的单位元，则称群胚 G 是单传递的。

命题 1. —— 设 G 为群胚。关系“存在 G 的箭连接 a 与 b”是 G 的顶点集上的等价关系，其等价类即 G 的轨道。

设 \(a\) 、\(b\) 、\(c\) 为 G 的顶点。我们有 \(e_a \in \mathrm{G}_{a,a}\) ；若 \(f \in \mathrm{G}_{a,b}\) ，则 \(f^{-1} \in \mathrm{G}_{b,a}\) ；若 \(f \in \mathrm{G}_{a,b}\) 且 \(g \in \mathrm{G}_{b,c}\) ，则有 \(fg \in \mathrm{G}_{a,c}\) 。这表明所指的关系是自反的、对称的且传递的；因此是等价关系。第二个断言由箭图连通分支的定义得出。

